% osglecture-adapters.dtx: structure-adapter definitions for osglecture.
% This file has no self-installing driver of its own; it is docstripped
% by the \generate block in osglecture.dtx (which must run in the same
% directory as this file). See osglecture.dtx for the license and the
% list of generated files.
%<*beamer>
\ProvidesExplFile
  {osglecture-adapter-beamer.def}
  {2026-08-21}
  {v0.13.0}
  {osglecture Beamer structure adapter}
\cs_new_eq:NN \__osglecture_adapter_beamer_lecture \lecture
\cs_new_protected:Npn \osglecture_adapter_lecture:nnn #1#2#3
  {
    \tl_if_blank:nTF {#1}
      { \__osglecture_adapter_beamer_lecture {#2} {#3} }
      { \__osglecture_adapter_beamer_lecture [#1] {#2} {#3} }
  }
% \ldeen{@2 trennt das Erfassen einer Anmerkung im Quelltext von ihrem
%   Ausgabeziel: ein Aufruf legt nur fest, \emph{dass} und \emph{mit
%   welchem Inhalt} eine Anmerkung existiert, nicht \emph{wohin} sie am
%   Ende geleitet wird. Aufrufstellen dürfen sich daher nicht auf ein
%   bestimmtes Ausgabeziel verlassen; welches Ziel @2 tatsächlich
%   bedient, wird allein adapterseitig entschieden und kann sich pro
%   Adapter oder im Zeitverlauf ändern, ohne dass Quelltext, der @2
%   benutzt, angepasst werden müsste. Für diesen Adapter ist @1 bereits
%   ein natives Ziel, das über mehrere Aufrufe pro Folie hinweg sammelt
%   statt zu überschreiben; @2 leitet lediglich unter eigenem Namen
%   dorthin weiter, damit das Ziel später geändert werden kann, ohne mit
%   @1s eigener Presenter-Notizen-Semantik zu kollidieren.}
%   {@2 decouples capturing a remark in the source from its eventual
%   output target: a call only records \emph{that} a remark exists and
%   \emph{with what content}, not \emph{where} it ends up. Call sites
%   must not assume any particular output target; which target @2
%   actually feeds is decided solely by the adapter and may differ
%   between adapters or change over time without requiring source that
%   uses @2 to be touched. For this adapter, @1 already provides a
%   native target that accumulates across multiple calls per slide
%   rather than overwriting; @2 merely forwards to it under its own
%   name, so the target can change later without colliding with @1's
%   own presenter-notes semantics.}
%   {\cs{note}}{\cs{Note}}
\cs_new_protected:Npn \osglecture_adapter_install:
  {
    \NewDocumentCommand \Note { m } { \note {##1} }
  }
\ExplSyntaxOff
%</beamer>
%<*book>
\ProvidesExplFile
  {osglecture-adapter-book.def}
  {2026-08-21}
  {v0.13.0}
  {osglecture book structure adapter}
\cs_new_eq:NN \__osglecture_adapter_book_chapter \chapter
\cs_new_eq:NN \__osglecture_adapter_book_frame \frame
\cs_new_eq:NN \__osglecture_adapter_book_endframe \endframe
\cs_new_protected:Npn \osglecture_adapter_lecture:nnn #1#2#3
  {
    \OsgLectureStandaloneTF
      { }
      {
        \tl_if_blank:nTF {#1}
          { \__osglecture_adapter_book_chapter {#2} }
          { \__osglecture_adapter_book_chapter [#1] {#2} }
      }
  }
% \ldeen{Die Zielbasisklasse kennt \env{frame}/\cs{frametitle} nicht mit
% Beamer-Semantik; @1 neutralisiert sie zu reinen Inhalts-Durchreichern
% (bzw. No-ops), damit gemeinsamer Quelltext mit \env{frame}-Umgebungen
% in der Langform ohne Sonderfall als gewöhnlicher Fließtext
% erscheint.}{The target base class has no \env{frame}/\cs{frametitle}
% with beamer semantics; @1 neutralizes them into plain content
% passthroughs (or no-ops), so shared source using \env{frame}
% environments appears as ordinary running text in the long form,
% without a special case.}{\cs{osglecture\_adapter\_install:}}
\cs_new_protected:Npn \osglecture_adapter_install:
  {
    \RenewDocumentEnvironment {frame} { D<> {all} O{} +b }
      {##3}
      { }
    \NewDocumentEnvironment {frame*} { D<> {all} O{} +b }
      {##3}
      { }
    \NewDocumentCommand \frametitle { D<> {all} O{##3} m } { }
    \NewDocumentCommand \framesubtitle { D<> {all} O{##3} m } { }
% \ldeen{Die Zielbasisklasse hat kein Presenter-Notizen-Ziel, an das
%   \cs{Note} weiterleiten könnte; @1 verschluckt die Anmerkung daher
%   spurlos, damit gemeinsamer Quelltext mit \cs{Note}-Aufrufen auch
%   hier ohne Sonderfall kompiliert.}{The target base class has no
%   presenter-notes target for \cs{Note} to forward to; @1 therefore
%   discards the remark silently, so shared source containing
%   \cs{Note} calls compiles here too without a special case.}
%   {\cs{Note}}
    \NewDocumentCommand \Note { m } { }
  }
\ExplSyntaxOff
%</book>
%<*ltx-talk>
\ProvidesExplFile
  {osglecture-adapter-ltx-talk.def}
  {2026-08-21}
  {v0.13.0}
  {osglecture ltx-talk structure adapter}
% \ldeen{@1 stellt bereits die gewünschten nativen Schnittstellen für
% \code{section}, \code{subsection}, \env{frame}, \env{frame*},
% \cs{frametitle} und \cs{framesubtitle} bereit; dieser Adapter lässt
% ihre Scanner deshalb bewusst unangetastet. Der Lehrkopf von @1
% übernimmt die von @2 erfassten Unit-Metadaten selbst.}{@1 already
% provides the desired native interfaces for \code{section},
% \code{subsection}, \env{frame}, \env{frame*}, \cs{frametitle}, and
% \cs{framesubtitle}; this adapter therefore deliberately leaves their
% scanners untouched. @1's own teaching header consumes the unit
% metadata recorded by @2.  The logical unit is also the presentation title;
% the profile-mapped course metadata supplies its subtitle.}{\cls*{ltx-talk}}{\pkg*{osglecture-structure}}
\cs_new_protected:Npn \osglecture_adapter_lecture:nnn #1#2#3
  {
    \tl_if_blank:nTF {#1}
      { \__osglecture_metadata_native_title {#2} }
      { \__osglecture_metadata_native_title [#1] {#2} }
  }
% \ldeen{@1 hat (noch) kein eigenes Presenter-Notizen-Ziel; @2
%   verschluckt die Anmerkung hier daher spurlos, statt eines von
%   fremder Semantik zu leihen. Das ist ein bewusst offener
%   Erweiterungspunkt: sobald @1 eine eigene Ausgabe für gesammelte
%   Anmerkungen bekommt, ersetzt deren Anbindung diesen No-op, ohne dass
%   sich an Aufrufstellen von @2 etwas ändern müsste.}{@1 has no
%   presenter-notes target of its own yet; @2 therefore discards the
%   remark silently here, rather than borrowing foreign semantics. This
%   is a deliberately open extension point: once @1 gains its own
%   output for collected remarks, wiring it in replaces this no-op
%   without requiring any change at call sites of @2.}
%   {\cls*{ltx-talk}}{\cs{Note}}
\cs_new_protected:Npn \osglecture_adapter_install:
  {
    \NewDocumentCommand \Note { m } { }
  }
\ExplSyntaxOff
%</ltx-talk>
%<*scrbook>
\ProvidesExplFile
  {osglecture-adapter-scrbook.def}
  {2026-08-21}
  {v0.13.0}
  {osglecture scrbook structure adapter}
\cs_new_eq:NN \__osglecture_adapter_scrbook_chapter \chapter
\cs_new_eq:NN \__osglecture_adapter_scrbook_frame \frame
\cs_new_eq:NN \__osglecture_adapter_scrbook_endframe \endframe
\cs_new_protected:Npn \osglecture_adapter_lecture:nnn #1#2#3
  {
    \OsgLectureStandaloneTF
      { }
      {
        \tl_if_blank:nTF {#1}
          { \__osglecture_adapter_scrbook_chapter {#2} }
          { \__osglecture_adapter_scrbook_chapter [#1] {#2} }
      }
  }
% \ldeen{Die Zielbasisklasse kennt \env{frame}/\cs{frametitle} nicht mit
% Beamer-Semantik; @1 neutralisiert sie zu reinen Inhalts-Durchreichern
% (bzw. No-ops), damit gemeinsamer Quelltext mit \env{frame}-Umgebungen
% in der Langform ohne Sonderfall als gewöhnlicher Fließtext
% erscheint.}{The target base class has no \env{frame}/\cs{frametitle}
% with beamer semantics; @1 neutralizes them into plain content
% passthroughs (or no-ops), so shared source using \env{frame}
% environments appears as ordinary running text in the long form,
% without a special case.}{\cs{osglecture\_adapter\_install:}}
\cs_new_protected:Npn \osglecture_adapter_install:
  {
    \RenewDocumentEnvironment {frame} { D<> {all} O{} +b }
      {##3}
      { }
    \NewDocumentEnvironment {frame*} { D<> {all} O{} +b }
      {##3}
      { }
    \NewDocumentCommand \frametitle { D<> {all} O{##3} m } { }
    \NewDocumentCommand \framesubtitle { D<> {all} O{##3} m } { }
% \ldeen{Die Zielbasisklasse hat kein Presenter-Notizen-Ziel, an das
%   \cs{Note} weiterleiten könnte; @1 verschluckt die Anmerkung daher
%   spurlos, damit gemeinsamer Quelltext mit \cs{Note}-Aufrufen auch
%   hier ohne Sonderfall kompiliert.}{The target base class has no
%   presenter-notes target for \cs{Note} to forward to; @1 therefore
%   discards the remark silently, so shared source containing
%   \cs{Note} calls compiles here too without a special case.}
%   {\cs{Note}}
    \NewDocumentCommand \Note { m } { }
  }
\ExplSyntaxOff
%</scrbook>
